\subsection{The remaining Dobrushin high-temperature
gap}\label{proof:positive_structure:the-remaining-dobrushin-high-temperature-gap}

For an Ising model, a Dobrushin condition such as
\(\max _i \sum_{j\ne i} \tanh (\beta  |J_{ij}|)<1\) controls boundary
effects; the underlying conditional-distribution method goes back to
\href{https://doi.org/10.1137/1113026}{Dobrushin (1968)}. Such
conditions yield spectral-gap and transportation estimates, for example
in \href{https://arxiv.org/abs/math/0611635}{Wu (2006)}. They do not
directly match the theorems above:

\begin{itemize}
\tightlist
\item
  Poincaré alone yields \(I_{\max}\ge v/(C n)\), retaining dimension
  dependence.
\item
  A rapidly mixing sampler need not provide an exact finite-degree,
  bounded-fanout source map. Truncating its history requires error
  control for all bounded low-degree observables, not just a per-site
  coupling bound.
\item
  Telescoping a sampler without bounding intermediate laws is invalid;
  the majority example shows this exactly.
\item
  If source fanout is \(L\), elementary path domination costs
  \(a^{-L}\). A polynomially large \(L\) in inverse approximation
  accuracy creates an exponential loss.
\item
  Total-variation approximation of the entire Gibbs law generally needs
  precision depending on \(n\). The block proof avoids that loss through
  low-degree interpolation; no analogous general Gibbs interpolation
  estimate has been proved here.
\end{itemize}

No general implication from Dobrushin uniqueness, spectral independence,
log-Sobolev inequalities, or efficient sampling alone to full polynomial
AA is established in this dossier. Such an implication remains an
identified research question, rather than a claimed consequence of
mixing.
